\section{Stage 1: Score-Function RL for Continuous Flows}
\label{sec:grpo}
The chain~(\ref{eq:chain}) permits a direct adaptation of GRPO \citep{shao2024deepseekmath} to continuous flow language models. At flow step $k$ the state is $s_{k-1}$ and the action is the next re-noised state $y_k$, sampled from the Gaussian transition $\mathcal N(y_k;\alpha_kD_\theta(s_{k-1}),\sigma_k^2I)$. During RL the readout is sampled rather than taken by $\arg\max$, $w_i\sim p_\theta(w_i\mid z_K)=\mathrm{Cat}(\mathrm{softmax}(\ell_i))$, so that the rollout has a joint density over its Gaussian transitions and its categorical readout; Stage 1 therefore optimises the objective~(\ref{eq:objective}) with the sampled readout in place of the $\arg\max$, and evaluation uses the $\arg\max$.

\paragraph{GRPO on the flow chain.}
We sample $G$ rollouts per prompt and form group-normalised advantages $A_g=\frac{r_g-\operatorname{mean}(r_{1:G})}{\operatorname{std}(r_{1:G})+\delta}$; the group mean, taken over all $G$ rollouts of the prompt including $g$, is the baseline. Each rollout batch is used for a single on-policy update, so no importance correction or clipping is required. We minimise the group-relative score-function loss
\begin{equation}
\mathcal L_{\mathrm{SF}}(\theta)=-\frac1G\sum_{g=1}^{G}A_g\Bigl[\sum_{k\in\mathcal K}\log\mathcal N\bigl(y^g_k;\alpha_kD_\theta(s^g_{k-1}),\sigma_k^2I\bigr)+\frac{1}{|\mathcal G|}\sum_{i\in\mathcal G}\log p_\theta\bigl(w^g_i\mid z^g_K\bigr)\Bigr],
\label{eq:grpo_loss}
\end{equation}
where $\mathcal K\subseteq\{2,\dots,K-1\}$ is a random subset of the transitions with a denoiser call, drawn per update to bound its cost, and the readout term is a per-token mean so that it does not grow with the response length. The transitions in $\mathcal K$ are recomputed from the stored states with the self-conditioning input $\hat x_{k-2}$ held fixed, and the readout term is recomputed from the stored $s_{K-1}$ through the final step and the decoder head. With the mean-centred advantage and no standard-deviation normalisation, the expected gradient at $\theta=\theta_{\mathrm{old}}$ is $(1-1/G)$ times the sum of two channels of the policy gradient of the sampled-readout objective: the flow-transition channel scaled by $|\mathcal K|/(K-2)$, with the dependence of later means on $\theta$ through the estimates omitted, and the readout channel scaled by $1/|\mathcal G|$ (Appendix~\ref{app:derivations}). Dividing by the within-group standard deviation trades this property for scale invariance across prompts \citep{liu2025drgrpo}.

% A clipped variant with per-step or per-position ratios supports off-policy reuse; ratio numerics in $L\times d$ dimensions are discussed in Appendix~\ref{app:derivations}.

% \paragraph{The credit-assignment design space.}
% Flow generation exposes two axes for assigning credit: flow step and token position. The coarsest scheme broadcasts one outcome-level advantage $A_g$ across all sampled transitions and generated positions. We generalize this scheme using flow-step advantages $A_{g,k}$, position advantages $A_{g,i}$, and dual-axis advantages $A_{g,k,i}$.Along the flow-step axis, process-aware credit assigns each transition an advantage derived from differences between successive intermediate values. We estimate these values using either a learned scalar critic or verifier scores on greedy intermediate readouts, the latter providing privileged task-specific credit. Along the position axis, each generated position receives a separate advantage that is shared across flow steps. These advantages come from a learned process reward model or token-aligned per-equation verifier labels measuring arithmetic validity, operand availability, and remaining-pool solvability. Combining both axes yields an advantage grid $A_{g,k,i}$ from changes in learned per-position values between successive states. We additionally evaluate sparse transition and grid credit, interpolations with the outcome-level advantage, and a clipped batch-reuse variant with per-position importance ratios. Together, these variants test whether finer credit assignment or ratio granularity can improve the score-function estimator.

\paragraph{Estimator analysis.}
The categorical-readout term in~(\ref{eq:grpo_loss}) has the usual token-level score, whereas each flow-transition term differentiates a Gaussian density over continuous embeddings. Because
$y_{k,i}=\alpha_kD_\theta(s_{k-1})_i+\sigma_k\varepsilon_{k,i}$, its score is
\begin{equation}
\nabla_\theta\log\mathcal N\!\left(
y_{k,i};
\alpha_kD_\theta(s_{k-1})_i,
\sigma_k^2I_d
\right)
=
\frac{\varepsilon_{k,i}^{\top}}{\sigma_k}
\alpha_k
\frac{\partial D_\theta(s_{k-1})_i}{\partial\theta}.
\label{eq:sf}
\end{equation}
Thus, the advantage scales the injected noise before it is propagated through the denoiser Jacobian. Unlike the categorical score of an AR action, an individual noise coordinate does not correspond to a discrete semantic decision. A coordinate contributes useful expected gradient only through its covariance with the advantage, yet it contributes sampling variance even when that covariance is negligible. Summing these terms across flow steps, generated positions, and embedding dimensions therefore produces an estimate whose signal comes from few coordinates and whose variance comes from all of them. A finer advantage, per step or per position, changes the weights of the terms but not the noise factor $\varepsilon_{k,i}/\sigma_k$ in each. The next stage removes that factor by differentiating the trajectory itself.